\section{Билет 48. Соотношение Майера. Уравнение внутренней энергии через постоянную адиабаты. Классическая теория теплоемкости идеального газа. Экспериментальная зависимость молярной теплоемкости при постоянном объеме от температуры.}

\begin{definition}
\textbf{Соотношение Майера} связывает молярные теплоёмкости идеального газа:
\begin{equation}
    C_p = C_V + R \quad \Rightarrow \quad C_V = \frac{R}{\gamma - 1}, \label{eq:mayer_gamma_48}
\end{equation}
где $\gamma = C_p/C_V$ --- показатель адиабаты. Подставляя $C_V$ из \eqref{eq:mayer_gamma_48} в формулу внутренней энергии идеального газа $U = \nu C_V T$, получаем выражение через макроскопические параметры:
\begin{equation}
    U = \nu \frac{R}{\gamma - 1} T = \frac{pV}{\gamma - 1}. \label{eq:U_gamma_48}
\end{equation}
Это соотношение широко используется в термодинамике для расчёта изменения внутренней энергии в любых процессах (не только изохорных), так как $U$ является функцией состояния.
\end{definition}

\begin{theorem}[Классическая теория теплоёмкости]
Согласно \textbf{закону равнораспределения энергии по степеням свободы} (классическая статистическая физика), при термодинамическом равновесии на каждую поступательную и вращательную степень свободы молекулы в среднем приходится кинетическая энергия $kT/2$.
Если молекула имеет $i$ степеней свободы (учитывая, что колебательные степени обладают удвоенной энергетической ёмкостью $kT$), средняя энергия одной молекулы равна $\langle \epsilon \rangle = \frac{i}{2} kT$. Тогда внутренняя энергия $\nu$ молей газа:
\begin{equation}
    U = \nu N_A \langle \epsilon \rangle = \nu \frac{i}{2} R T.
\end{equation}
Молярная теплоёмкость при постоянном объёме определяется как $C_V = \frac{1}{\nu} \frac{dU}{dT}$. Отсюда получаем \textbf{классическую формулу}:
\begin{equation}
    C_V = \frac{i}{2} R. \label{eq:classical_Cv}
\end{equation}
Согласно классической теории, $C_V$ --- постоянная величина, не зависящая от температуры, давления и объёма газа.
\end{theorem}

\begin{remark}
\textbf{Экспериментальная зависимость $C_V(T)$.} Опыт показывает, что формула \eqref{eq:classical_Cv} справедлива лишь в ограниченном диапазоне температур. На самом деле $C_V$ зависит от $T$ ступенчато:
\begin{enumerate}
    \item \textbf{Низкие температуры} ($T \to 0$): Тепловой энергии $kT$ недостаточно для возбуждения вращательных и колебательных степеней свободы. Они ``замораживаются''. Газ ведёт себя как одноатомный ($C_V \approx \frac{3}{2}R$).
    \item \textbf{Умеренные температуры} ($T \sim 300$ К): Вращательные степени свободы ``размораживаются''. Для двухатомных газов $i$ возрастает с 3 до 5, и $C_V \approx \frac{5}{2}R$.
    \item \textbf{Высокие температуры} ($T > 1000$ К): Возбуждаются колебательные степени свободы. $C_V$ снова возрастает, стремясь к $\frac{7}{2}R$ (для двухатомного газа).
\end{enumerate}
Эта зависимость объясняется только \textbf{квантовой теорией}: энергия вращения и колебаний квантуется, и для перехода на следующий энергетический уровень требуется преодолеть определённый энергетический порог. Классическая теория является предельным случаем квантовой при высоких температурах ($kT \gg \Delta \epsilon$).
\end{remark}
